\section{Требования к системе навигации и допущения мезоскопической симуляции}

Для реализации концепции глобального маршрутизирующего узла, способного обслуживать сотни тысяч транспортных средств в реальном времени, необходимо сформировать жесткие технические и архитектурные требования. 

\subsection{Базовые требования к высокопроизводительному роутингу}

Пропускная способность маршрутизатора является критическим узким местом во всех существующих транспортных симуляторах. Для достижения целевого показателя пропускной способности (свыше 5000 запросов в секунду, QPS) в разрабатываемом комплексе были заложены следующие архитектурные принципы:

\begin{enumerate}
    \item \textbf{Абсолютная изоляция состояний (Stateless-архитектура).} Вычислительный контур алгоритма маршрутизации не должен зависеть от глобального состояния симуляции. В горячем цикле поиска путей (цикле алгоритма A*) полностью исключается использование мьютексов (mutex), семафоров и любых других механизмов межпоточной синхронизации. Каждый поток маршрутизатора работает исключительно с локальной памятью и атомарными операциями чтения предрассчитанных констант.
    \item \textbf{Устранение накладных расходов на ввод-вывод (Zero-copy).} Традиционный подход, при котором граф считывается из базы данных или дискового хранилища в оперативную память при старте приложения, требует огромных объемов RAM и увеличивает время старта системы. В разрабатываемой системе архитектура графа запекается в плоские бинарные файлы (CSR-структуры), доступ к которым осуществляется через механизм отображения файлов в память (системный вызов \texttt{mmap}). Это передает управление кэшированием непосредственно ядру операционной системы Linux (Page Cache) и обеспечивает мгновенный (zero-copy) доступ к геоданным без выделения дополнительной памяти (zero-allocation) внутри приложения.
    \item \textbf{Поддержка динамических весов без инвалидации индексов.} Многие популярные системы маршрутизации (например, использующие метод Contraction Hierarchies) достигают высокой скорости поиска за счет добавления в граф так называемых шорткатов (shortcuts). Любое изменение скорости движения на реальной дороге делает весь предрассчитанный граф недействительным. Разрабатываемая система должна использовать эвристические подходы (такие как ALT — A* search, Landmarks, and Triangle inequality), которые математически сохраняют свою допустимость даже при динамическом увеличении весов ребер из-за возникших заторов, не требуя перестроения вспомогательных структур.
\end{enumerate}

\subsection{Допущения для мезоскопического моделирования}

Создание физически достоверной модели движения в масштабах всего города является вычислительно неподъемной задачей, если моделировать поведение каждой машины на микроуровне. Для обеспечения баланса между производительностью и достоверностью (наблюдаемостью образования заторов), в рамках работы применяется \textbf{мезоскопический подход} к моделированию, базирующийся на следующих осознанных допущениях:

\begin{itemize}
    \item \textbf{Отказ от декартовой координатной сетки.} Автомобили в симуляции не имеют точных физических координат $(x, y)$ и угла поворота кузова. Пространство жестко дискретизируется до гранулированных направленных сегментов дорожного графа (ребер). Движение агента рассматривается как одномерное перемещение вдоль конкретного ребра от его начала (позиция $S = 0$) до конца (позиция $S = Length$).
    \item \textbf{Отсутствие микроскопической физики взаимодействия.} Из модели полностью исключаются такие сложные поведенческие паттерны водителей, как перестроение между полосами, поддержание безопасной дистанции следования (модель car-following), обгоны и микрозадержки реакции человека. Агент в симуляции — это «частица» потока идеального водителя.
    \item \textbf{Дискретизация времени.} Время в моделировании разделяется на фиксированные кванты. На микроуровне (внутри кинематического движка) расчеты положения производятся с заданным шагом $dt$ (например, 10 или 50 миллисекунд). Однако на макроскопическом уровне оценки плотности потока и формирования заторов время квантуется на более крупные 5-минутные интервалы (Volume Buckets).
    \item \textbf{Гомогенность агентов.} Для упрощения расчетов предполагается, что все легковые автомобили имеют усредненную эффективную длину (например, 5 или 7 метров, с учетом необходимого межмашинного интервала в заторе). Это позволяет использовать жесткую метрику пространственной вместимости сегмента дороги.
\end{itemize}

Данный набор допущений отсекает колоссальный пласт микроскопических вычислений, позволяя сместить фокус вычислительной мощности CPU на точный расчет гидродинамических волн заторов (Spillback), что является ключевым индикатором здоровья транспортной сети на макроуровне.
